\section{Modelos de Difusión Variacional: un marco unificado}

\subsection{Forma general}

Los Modelos de Difusión Variacional definen una familia genérica de modelos de difusión. La forma general de la distribución de salto es:

$$
q(\mathbf{x}_t \mid \mathbf{x}_0) = \mathcal{N}\!\left(\mathbf{x}_t;\, \alpha_t \mathbf{x}_0,\, \sigma_t^2 \mathbf{I}\right)
$$

donde $\alpha_t$ y $\sigma_t$ son funciones diferenciables del tiempo $t$, que satisfacen:

\begin{enumerate}
    \item $\alpha_0 = 1$, $\sigma_0 = 0$ (en el momento inicial se trata de los datos originales);
    \item $\alpha_T \approx 0$, $\sigma_T \approx 1$ (en el momento final es aproximadamente una gaussiana estándar);
    \item \textbf{La relación señal-ruido SNR$(t) = \alpha_t^2/\sigma_t^2$ decrece monótonamente con respecto a $t$}.
\end{enumerate}

\begin{importantbox}{Significado de la disminución monótona del SNR}
El hecho de que el SNR decrezca monótonamente significa que, a medida que avanza el proceso hacia adelante, el componente de señal ($\alpha_t \mathbf{x}_0$) se debilita continuamente en relación con el componente de ruido ($\sigma_t \boldsymbol{\epsilon}$). Esta es la esencia del "añadido de ruido" en los modelos de difusión: la información se va sumergiendo progresivamente en el ruido, hasta que la distribución de los datos degenera en una distribución de puro ruido.
\end{importantbox}

\subsection{SMLD y DDPM como casos particulares}

\begin{table}[H]
\centering
\begin{tabular}{lccc}
\toprule
\textbf{Modelo} & $\alpha_t$ & $\sigma_t$ & \textbf{Tipo de SDE} \\
\midrule
SMLD / NCSN & $1$ & $\sigma(t)$ & VE-SDE \\
DDPM & $e^{-\frac{1}{2}\int_0^t \beta(s)\,\mathrm{d}s}$ & $\sqrt{1 - e^{-\int_0^t \beta(s)\,\mathrm{d}s}}$ & VP-SDE \\
\bottomrule
\end{tabular}
\caption{Parámetros de SMLD y DDPM bajo el marco de los Modelos de Difusión Variacional}
\end{table}

\begin{knowledgebox}{Apertura del espacio de diseño}
VE-SDE y VP-SDE son solo dos ejemplos concretos. Cualquier par $(\alpha_t, \sigma_t)$ que cumpla la condición de disminución monótona del SNR define un modelo de difusión válido. Esto ofrece a los investigadores un enorme espacio de diseño: pueden optimizar la calidad de generación, la velocidad de muestreo u otros indicadores eligiendo diferentes agendas de ruido. Por ejemplo, EDM (Elucidating the Design Space of Diffusion-Based Generative Models, Karras et al. 2022) exploró sistemáticamente este espacio de diseño.
\end{knowledgebox}

\subsection{De SDE a distribución de salto}

Dadas $f(\mathbf{x}, t)$ y $g(t)$ del SDE hacia adelante, se pueden deducir los correspondientes $\alpha_t$ y $\sigma_t$:

$$
\alpha_t = e^{\int_0^t f(s)\,\mathrm{d}s}, \qquad \sigma_t^2 = \alpha_t^2 \int_0^t \frac{g^2(s)}{\alpha_s^2}\,\mathrm{d}s
$$

A la inversa, dados $\alpha_t$ y $\sigma_t$, también se pueden recuperar $f(t)$ y $g(t)$:

$$
f(t) = \frac{\mathrm{d} \ln \alpha_t}{\mathrm{d}t}, \qquad g^2(t) = \frac{\mathrm{d}\sigma_t^2}{\mathrm{d}t} - 2\frac{\mathrm{d} \ln \alpha_t}{\mathrm{d}t}\sigma_t^2
$$

Esta conversión bidireccional implica que existen \textbf{tres descripciones equivalentes} —distribución de transferencia, distribución de salto y SDE— que pueden intercambiarse libremente.

\begin{warningbox}{Advertencia sobre confusión de símbolos}
Diferentes artículos utilizan convenciones de simbolismo distintas. Por ejemplo, el $\sigma$ de Song et al. se refiere a la agenda de ruido en VE-SDE (es decir, el $\sigma_i$ de SMLD), mientras que el $\sigma_t$ en los Modelos de Difusión Variacionales es la desviación estándar de la distribución de salto. En un mismo desarrollo pueden aparecer simultáneamente ambos tipos de $\sigma$, y deben distinguirse según el contexto. El profesor Minhyuk Sung también advirtió especialmente sobre este punto durante la clase.
\end{warningbox}

\subsection{Resumen de la sección}

Los Modelos de Difusión Variacional proporcionan un marco general mediante la parametrización $(\alpha_t, \sigma_t)$; SMLD y DDPM son sus casos particulares. Existe una correspondencia bidireccional exacta entre las tres descripciones (distribución de transferencia, distribución de salto y SDE).
